\section{Actual prime chains and a medium-prime central barrier}
\label{sec-ng18-prime-chain}

For the actual endpoint $c=2^e3^f$, with integers $e,f\geq2$, put
\[
 M=2^e,\quad N=3^f,\quad A=M/2,\quad B=N/3,
 \qquad Q=\operatorname{rad}(c-1).
\]
Retain both complete premises
\[
 \max(M,N)<2\min(M,N)-1,\qquad AB>Q.                 \tag{NG0}
\]
The previously proved DP no-proper-face calculation identifies the
original FCRT/SCRT/PBT optimum as $D+\Gamma$, where
\[
 D=\log(AB/Q)>0,\qquad
 \Gamma=\min_{H\mid Q}\operatorname{dist}
       \bigl(\log H,[\log(Q/B),\log A]\bigr).
\]
All divisors here are divisors of the complete actual radical. Each
prime is assigned once; in the original endpoint this assigns its
entire prime-power contribution. The following results concern that
same optimum.

\paragraph{NG1: the exact largest divisor gap.}
Let $Q>1$ be squarefree, with primes $p_1<\cdots<p_s$, and define
\[
 R_i=\prod_{j<i}p_j,\qquad
 \Lambda(Q)=\max_{1\leq i\leq s}\frac{p_i}{R_i}.
\]
The largest difference of consecutive logarithms of divisors of $Q$
is exactly $\log\Lambda(Q)$.

\emph{Proof.} When $p_i$ is added, the previous divisor-log set $S$
is replaced by $S\cup(S+a)$, where $a=\log p_i$ and $S$ has endpoints
$0,T=\log R_i$. If $a>T$, the spans are disjoint and the largest gap
is the maximum of the old gap and $a-T$. If $a\leq T$, the spans
overlap. Every gap in the union is contained in a gap of at least one
of the two sets, within that set's span, so it is at most the old
maximum. Induction, starting with $\{0,\log p_1\}$, proves the upper
bound.
For the lower bound, if $p_i>R_i$, a divisor of the final $Q$ smaller
than $p_i$ cannot contain $p_i$ or any later prime. It divides $R_i$
and is at most $R_i$. Thus $R_i,p_i$ are consecutive divisors of the
final integer. Since $p_1>1=R_1$, the largest of these positive gaps
attains $\log\Lambda(Q)$. $\square$

Once the full prime support is supplied, this is a test with one
step per prime. It makes no assertion about the cost of factoring $Q$.

\paragraph{NG2: a sufficient arithmetic class.}
For any real $A,B\geq1$ with $AB>Q$, the same definitions give
\[
 \Gamma\leq\frac{\max(0,\log\Lambda(Q)-D)}2
 \leq\frac{\max(0,\log P^+(Q)-D)}2.                 \tag{NG1}
\]
In particular,
\[
 p_i\leq (AB/Q)R_i\quad\hbox{for every }i             \tag{NG2}
\]
implies $\Gamma=0$.

\emph{Proof.} If $\Gamma>0$, the central closed interval contains
neither $0$ nor $\log Q$, since both are divisor logs. Because
$A,B\geq1$, it therefore lies strictly inside $[0,\log Q]$, in one
open gap $(u,v)$ of the full divisor-log set. Its length is $D$.
The two endpoint distances sum to $v-u-D$, so their minimum is at
most $(\log\Lambda(Q)-D)/2$. If $\Gamma=0$, the bound is immediate.
Finally $R_i\geq1$ gives $\Lambda(Q)\leq P^+(Q)$, and (NG2) is
exactly $\log\Lambda(Q)\leq D$. $\square$

This proves a full arithmetic subclass, but not its membership for
all endpoints. A large gap near an end of the divisor set can also
make (NG2) fail when the required central partition succeeds.

\paragraph{NG3: a certified nonsquare central success.}
Take $e=41,f=26$. The complete factorization is
\[
 c-1=7^2\cdot439\cdot857\cdot2729\cdot292183\cdot380261663.
                                                               \tag{NG3}
\]
All six factors are prime, with the deterministic certificates
described below. Both premises (NG0) hold. The actual divisors
\[
 H=2729\cdot380261663=1037734078327,\qquad
 J=7\cdot439\cdot857\cdot292183
\]
satisfy $HJ=Q$, $H\leq2^{40}=A$ and $J\leq3^{25}=B$ by integer
comparison. Hence $\Gamma=0$, while
$D=\log(7c/(6(c-1)))>0$. Since $e$ is odd, this endpoint is outside
the earlier DP square family. Its global value $\Lambda(Q)=439/7$
is larger than $AB/Q$. Thus (NG2) does not certify every successful
central partition. This is one certified nonsquare endpoint, not an
assertion of infinitely many such endpoints.

\paragraph{NG4: a complete medium-prime bracket.}
Let $Q=qR$ be squarefree, with $q$ prime and $p<r$ the two smallest
prime factors of $R$. Suppose $A,B\geq1$, $AB>Q$, and
\[
 prq\leq R<r^2q,\qquad rq>\max(A,B).                \tag{NG4}
\]
Then $R/r,rq$ are consecutive divisors of the full $Q$, and
\[
 \Gamma=\log\frac{rq}{\max(A,B)}>0,\qquad
 D+\Gamma=\log\frac{\min(A,B)}{R/r}.                \tag{NG5}
\]

\emph{Proof.} A divisor of $R$ below $r$ is either $1$ or $p$,
because $R$ is squarefree. If $H\mid Q$ contains $q$ and $H<rq$,
write $H=qd$ with $d\mid R,d<r$. Then $H\leq pq\leq R/r$.
If $H$ omits $q$ and $H>R/r$, then $R/H<r$, so $H=R$ or $R/p$;
both are at least $rq$. The endpoints themselves divide $Q$, are
strictly ordered, and have product $Q$. Since $rq>A,B$, the entire
central interval lies strictly between their logarithms. Its two
distances are $\log(rq/B)$ and $\log(rq/A)$, proving both formulas.
$\square$

For an actual instance take $e=123,f=78$, the cube of the preceding
endpoint. Put
\[
 \begin{split}
 q&=257077829249200341434761489003997119,\\
 R&=7\cdot31\cdot439\cdot857\cdot919\cdot2729
      \cdot292183\cdot380261663\cdot4266021703.
 \end{split}
\]
The complete factorization is $c-1=7qR$, hence $Q=qR$ without an
unfactored remainder. Here $p=7,r=31$; both (NG0) and (NG4) hold.
More strongly, every prime of $Q$ is at most $q$ and
\[
 20q<\min(A,B),                                      \tag{NG6}
\]
where the exact integers are
\[
 \begin{aligned}
 A&=5316911983139663491615228241121378304,\\
 B&=5474401089420219382077155933569751763,\\
 R/31&=3130566812026385226903752016561167657,\\
 31q&=7969412706725210584477606159123910689.
 \end{aligned}
\]
Since $B>A$, the exact conclusions are
\[
 D=\log\frac{7c}{6(c-1)},\qquad
 \Gamma=\log\frac{31q}{B}>2D>0,\qquad
 D<\log(6/5)<1.                                      \tag{NG7}
\]
No decimal logarithm is needed: the two strict comparisons are
\[
 31q[6(c-1)]^2>B(7c)^2,\qquad
 35c<36(c-1).
\]
The exact positive balance margin
$2\min(M,N)-1-\max(M,N)$ is
\[
 4844444664297995820229445163776257926.
\]
The verifier checks these comparisons and all $1024$ actual divisors.
The proof of (NG5) separately establishes completeness of the bracket.

Consequently (NG0), nonsquareness, $0<D<1$, and even
$P^+(Q)<\min(A,B)/20$ do not imply $\Gamma=0$ or $\Gamma\leq D$.
This excludes precisely those strengthened statements. It supplies
neither a fixed-$\epsilon$ infinite counterexample nor a contradiction
to a bound with an arbitrary additive $C_\epsilon$. It does not
disprove ABC or another certificate with different, proved rules.
Within the original rules, (NG5) is the exact optimum.

\paragraph{NG5: deterministic primality and complete finite evidence.}
The standard-library verifier
\texttt{next\_replay\_medium\_\allowbreak{}partition.py} checks
47 prime nodes and 134 witnesses: base $2$, 45 full predecessor
nodes and one full Gaussian successor node. It uses no probable-prime
predicate or factorization routine. For completeness, both criteria
have elementary proofs.

For a predecessor node $m>2$, factor $m-1$ completely into previously
proved primes. For each distinct $r\mid m-1$, suppose an integer
$a_r$ satisfies
\[
 a_r^{m-1}\equiv1\pmod m,\qquad
 \gcd(a_r^{(m-1)/r}-1,m)=1.
\]
For every prime $\ell\mid m$, the order of $a_r$ modulo $\ell$
divides $\ell-1$ and $m-1$ and contains the full $r$-part of $m-1$.
It follows that $m-1\mid\ell-1$, hence $\ell\geq m$ and $\ell=m$.
Thus $m$ is prime. A different witness for each $r$ is permitted.

For a successor node let $m>2$ be odd and factor $m+1$ completely
into previously proved smaller primes. For every $r\mid m+1$,
suppose $z_r=a_r+b_ri$ in $(\mathbb Z/m\mathbb Z)[i]$ satisfies
\[
 a_r^2+b_r^2\equiv1\pmod m,\qquad z_r^{m+1}=1,\qquad
 \gcd\bigl(N(z_r^{(m+1)/r}-1),m\bigr)=1.             \tag{NG8}
\]
For any prime $\ell\mid m$, the algebra $\mathbb F_\ell[i]$ is
either $\mathbb F_\ell\times\mathbb F_\ell$ or
$\mathbb F_{\ell^2}$, since $\ell$ is odd. Its norm-one group has
order $\ell-1$ in the split case, where its elements are
$(u,u^{-1})$, and $\ell+1$ in the field case, by the surjective
finite-field norm. The last condition in (NG8) makes
$z_r^{(m+1)/r}-1$ a unit modulo $\ell$, so the order of $z_r$ contains
the full $r$-part of $m+1$. Therefore $m+1$ divides that norm-one
group order for every $\ell\mid m$, giving $\ell\geq m$ and hence
primality. No Jacobi-symbol test is being substituted for this proof.

The sole successor node is
\[
 \begin{split}
 m&=13098839766085821942054493478243,\\
 m+1&=2^2\cdot3^5\cdot7\cdot23\cdot1427\cdot10691
       \cdot40507\cdot135446661817793.
 \end{split}
\]
All factors are earlier proved nodes. The largest endpoint prime
has $q-1=2\cdot3\cdot3271\cdot m$, a full predecessor node.
The Gaussian operations are literal integer-pair operations modulo
$m$, and every supplied witness is checked.

The verifier also checks both complete factorizations of $c-1$,
primitivity, balance, nonsquareness, $D>0$, all $64$ and $1024$
divisors, the maximum-gap formula, the central memberships and the
strict rational comparisons. Generation and read-only
\texttt{--check} actually passed; independent root and peer executions
also passed. The canonical output SHA256 is
\texttt{26d04f39748d763f\allowbreak{}5139794465afb780\allowbreak{}58f18474567b629f\allowbreak{}6937abcae2ec9af2}.
These finite certificates and the complete elementary proofs above
are separate from any Lean theorem, prime-value infinitude, or
uniform ABC estimate.
